% Learning InitechBase, Chapter 5: Changing Your Data (AN-16).  Every panel
% replayed (examples/l5-*.dot); sources cmd/mutate.c (APPEND BLANK,
% REPLACE, DELETE, RECALL, PACK), cmd/set.c (DELETED).
\chapter{Changing Your Data}
\Gates{test-pubs-examples, test-interp-replace, test-dbf-mutate}
Records change: a new employee joins, an extension moves, salaries rise, an
employee leaves.  In this lesson you add a record, change records one at a time
and many at once, and mark records for deletion and remove them.

\begin{ChapterCovers}
\item Adding a record with \cmd{APPEND BLANK} and \cmd{REPLACE}
\item Changing one record
\item Changing many records at once
\item Marking records with \cmd{DELETE}, and \cmd{RECALL}
\item SET DELETED
\item Removing records with \cmd{PACK}
\end{ChapterCovers}

\ATEHead{Preparing for This Chapter}
The lesson changes \file{STAFF.DBF}.  Start InitechBase with a fresh copy of the
sample files, and copy them again whenever you want to begin again.

\ATEHead{Adding a Record}
\cmd{APPEND BLANK} adds an empty record at the end of the database and makes it
the current record.  \cmd{REPLACE} then fills in its fields: after
\cmd{REPLACE} comes each field, the word \kw{WITH} and the value.  Character
values are in quotation marks; a date is converted with \fn{CTOD()}; a logical
value is \texttt{.T.} or \texttt{.F.}.
\begin{Steps}
\item Type \type{USE STAFF INDEX STAFFNAM, STAFFNO}.
\item Type \type{APPEND BLANK}.
\item Type \type{REPLACE EMPNO WITH '1014', LASTNAME WITH 'HARRIS', FIRSTNAME
WITH 'JOAN'}.
\item Type \type{REPLACE DEPT WITH 'ACCT', EXT WITH '1432', HIRED WITH
CTOD('03/02/92')}.
\item Type \type{REPLACE SALARY WITH 31500, EXEMPT WITH .F.}.
\item Type \type{DISPLAY EMPNO, LASTNAME, FIRSTNAME, DEPT, HIRED, SALARY}, then
\type{? RECCOUNT()}.
\end{Steps}
\DotExample{examples/l5-append}
The new record is written to the disk as soon as it is added and changed.  The
index files open with the database are kept up to date, so the new employee is
in alphabetical order at once:
\DotExample{examples/l5-indexed}

\ATEHead{Changing One Record}
Find the record, then \cmd{REPLACE} the field.
\begin{Steps}[resume]
\item Type \type{USE STAFF}, then \type{GO 8}.
\item Type \type{REPLACE EXT WITH '1404'}, then \type{DISPLAY LASTNAME, EXT}.
\end{Steps}
\DotExample{examples/l5-replace}

\ATEHead{Changing Many Records}
With a \kw{FOR} clause, \cmd{REPLACE} changes every record that meets the
condition.  The value is worked out for each record in turn, so a field may be
used to compute its own new value.  Here everyone in Customer Relations
receives a raise of four per cent:
\begin{Steps}[resume]
\item Type \type{USE STAFF}.
\item Type \type{REPLACE SALARY WITH SALARY * 1.04 FOR DEPT = 'CUST'}.
\item Type \type{LIST LASTNAME, SALARY FOR DEPT = 'CUST'}.
\end{Steps}
\DotExample{examples/l5-raise}
\begin{AteWarning}
\cmd{REPLACE} with \kw{FOR} changes every record that meets the condition, and
there is no command that undoes it.  List the records with the same \kw{FOR}
clause before you replace them.
\end{AteWarning}

\ATEHead{Marking Records for Deletion}
\cmd{DELETE} does not remove a record: it marks it.  A marked record is listed
with an asterisk and \fn{DELETED()} is true for it.
\begin{Steps}[resume]
\item Type \type{USE STAFF}.
\item Type \type{LOCATE FOR LASTNAME = 'FENWICK'}, then \type{DELETE}.
\item Type \type{LIST LASTNAME, DEPT FOR DELETED()}.
\end{Steps}
\DotExample{examples/l5-delete}
Record 13, Robert Chen, was already marked on the diskette.  While SET DELETED
is ON, InitechBase passes over the marked records:
\DotExample{examples/l5-setdeleted}
\cmd{RECALL} removes the mark from the current record:
\DotExample{examples/l5-recall}

\ATEHead{Removing Records}
\cmd{PACK} removes the marked records for good and renumbers the others.
\begin{Steps}[resume]
\item Type \type{USE STAFF}, then \type{? RECCOUNT()}.
\item Type \type{PACK}, then \type{? RECCOUNT()}.
\item Type \type{GO BOTTOM}, then \type{? RECNO(), LASTNAME}.
\end{Steps}
\DotExample{examples/l5-pack}
\begin{AteWarning}
Records removed by \cmd{PACK} cannot be recalled.  \cmd{ZAP} removes all the
records of a database at once; see \doctitle{Using InitechBase}, Chapter~5.
\end{AteWarning}

\ATEHead{Where Do You Go From Here?}
You can keep a database up to date.  In Chapter~6 you let InitechBase calculate
for you.

\ATEHead{Summary}
\begin{enumerate}
\item \cmd{APPEND BLANK} adds an empty record; \cmd{REPLACE} \xp{field}
  \kw{WITH} \xp{value} fills it in.
\item \cmd{REPLACE} with \kw{FOR} changes every record that meets a condition.
\item \cmd{DELETE} marks a record; \cmd{RECALL} removes the mark; SET DELETED ON
  hides marked records.
\item \cmd{PACK} removes the marked records.
\item Changes are written to the disk at once, and the open index files are kept
  up to date.
\end{enumerate}
